% Français 9115, batch 5 — Gallica views 81–100.
% Pass: Opus 5 (claude-opus-5), 2026-09-12 — first pass, unchecked against
% the views by a human.
%
% What the twenty views hold. Two things, and the seam between them falls
% inside the batch. Views 81 to 88 finish the fair copy of the « Mémoire sur
% la courbure des surfaces » that batch 3 opened at view 43; view 89 is the
% plate of its four figures, drawn in ink and carrying no text, which the
% memoir has been calling « fig. 1 » to « fig. 4 » since view 73. Then the
% fair copy stops. View 90 is blank, view 91 carries nothing but the title
% « Moyennes courbures » in a large hand, view 92 is its blank verso — and
% from view 93 to the end of the batch come working drafts of the same
% subject, in a fast, heavily corrected hand, one written leaf at a time with
% its verso left blank: 93, 95, 97 and 100 written, 94, 96, 98 and 99 blank.
% Ten views of the twenty carry text.
%
% The argument of the fair copy's last pages. Batch 4 broke off inside N° 8,
% at the surfaces whose principal radii lie opposite ways.
%
%   v. 81–82   That for every figure between the cylindrical form and the one
%              whose principal curvatures are equal and of contrary signs the
%              angles of the development do not change, the interval $pq$
%              being by construction a quarter of $AC$; so the directrix is
%              always the cylinder's, and the figures differ only in how much
%              of the directrix lies above the tangent plane. Passing the one
%              into the other, the « surface des distances » keeps its form
%              while the « surface des distances moyennes » wastes away and
%              vanishes when the lines of distance on the two faces of the
%              tangent plane are equal.
%   v. 83–84   That the law of the distribution of curvature is unaltered by
%              the change: the sum of the lines of distance in two
%              perpendicular normal planes is independent of the choice of
%              those planes, the line at $45^\circ$ is the mean, and four
%              directions give equal lines of distance — two only for the
%              principal ones. It is enough, to carry the earlier proofs
%              over, to count the increase above $AC$ as a decrease.
%              Then N° 9, and the turn to mechanics: without the law of the
%              distribution and the mean curvature bound up with it, no
%              question where curvature acts as a dynamical quantity could be
%              treated in its generality.
%   v. 85–86   Forces proportional to the curvature are spread over the whole
%              contour of a point, infinite in number, each infinitely small
%              and proportional to the line of distance in the plane of its
%              direction. The action exerted is therefore independent of how
%              they are distributed, and equal to what a uniform distribution
%              would give — that is, the action the point would feel on the
%              sphere of mean curvature. Whence the proposition the memoir
%              has been building towards: the action is proportional to the
%              « surface des distances moyennes ».
%   v. 87–88   Where the composition of forces gives convenient formulas, the
%              « surfaces des distances », taken as solids wedged between the
%              surface and its tangent plane, give the picture. A warning not
%              to confound the figure of a surface with the curvature taken
%              as a dynamical quantity; then the applications, cut short — of
%              vaults, where the small number of forms has hidden from the
%              geometers the need to express the resistances born of
%              curvature in general, and dômes against berceaux; and last the
%              example of two surfaces of unlike figure with equal extent and
%              the same mean curvature at every point, the circumscribed
%              cylinder and the quarter of the sphere of double radius.
%   v. 89      The plate. Four figures, each a parallelogram $ABDC$ crossed by
%              the broken line, lettered as the text has lettered them.
%
% The drafts. View 93 takes up the vaults again — a dome as a solid of
% revolution, its zones matched to cones and then to cylindrical strips of
% the same curvature, with two segments $A$——$B$ and $A'$——$B'$ drawn on the
% leaf and a brace tying the lines together; the leaf is torn across and the
% foot of it is gone. Views 95, 97 and 100 are drafts of the passage where
% the memoir meets the equation of elastic surfaces: the space run through by
% the points of a solid whose form an external cause has changed (95); that
% the displacement of the surface's points gives that of the points in the
% depth, so that thickness and natural elasticity enter only as a constant
% coefficient of two factors multiplying the equation found for the surface
% (97); and the two distances, before and after the change of form, bound
% together by the law that constitutes the notion of mean curvatures (100).
% View 95 carries « Commencement de la p. A », an instruction of the author
% to herself, and the letter « E » at the head; view 97 carries in its left
% margin the addition « cette uniformité elle même résulte de
% l'invariabilité de l'épaisseur ». These leaves are the working state of
% what the fair copy says smoothly, and they are transcribed with their
% deletions.
%
% The numbering on the leaves, and why no \folio{} is written. The ink
% series of batches 1 to 4 runs on over the fair copy, on the rectos: 40, 41,
% 42, 43, 44 at views 81, 83, 85, 87, 89. Beside the first three stands the
% author's own pagination, odd and struck through, of which 39 was read at
% view 83 and the others are struck past reading; at view 87 a single « 43 »
% stands, unstruck. The even pagination is in the top left corner of each
% verso, unstruck: 38, 40, 42, 44 at views 82, 84, 86, 88. The draft leaves
% carry the same ink series and nothing else: 45 at view 91, 46 at view 93,
% 47 at view 95 (with « E » beside it), 48 at view 97, 49 at view 99 — that
% last on the blank side of the leaf whose written side is view 100. Every
% number was read on the view; none is inferred. All are in ink: no pencilled
% foliation of the Bibliothèque was found anywhere in the batch, and no
% \folio{} is therefore written, as in batches 1 to 4.
%
% Notation. The letters are those of the plate: $A$, $B$, $C$, $D$ for the
% parallelogram, $p$, $q$, $h$, $g$, $m$, $n$ and their primes for the points
% of the development. The broken line is written « $Bqh$ \&c » at views 81
% and 82. One inconsistency is left standing at view 84, where the four
% points of the development are named $v$, $t$, $v'$, $t'$ in one paragraph
% and $v$, $u$, $v'$, $u'$ in the next; the plate has $v$, $u$, $v'$, $u'$.
% « qui agissant sur un point » at view 84 is likewise hers.
%
% Her hand, and what the leaves do to it. The fair copy is the regular hand
% of batches 3 and 4; the drafts are another matter — a fast hand, the
% terminal flourishes longer, words overwritten rather than struck, and
% additions squeezed between the lines and into the margins. Her own
% underlining is set \emph{}; \uncertain{} renders as an underline too, so a
% reader of the PDF should take this header's word for which is which. At
% views 86 and 88 the right edge of the leaf passes under the mount and the
% last letters of several lines are hidden: the words are given whole where
% the sense leaves no room for doubt and the fact is noted at the spot.
%
% What could not be read. On the fair copy, almost nothing: the reading
% « il en est » at view 83, the struck word under « chacune » at view 85, and
% the word overwritten by « avec » at view 87. On the drafts, a good deal,
% and all of it marked: the three struck words after « N° » at view 95; the
% four lines added afterwards in a very fine hand at the head of view 97,
% which are read only in part; the struck clusters at views 97 and 100; and
% the last line of view 100, reduced to a word and to ink blots. Nothing was
% guessed to fill a gap.
%
% Nothing here was checked against any published transcription.
\documentclass[11pt,a4paper]{article}
% The preamble shared by every transcription of the mathematicians' manuscripts in Gallica.
%
% It rests on one idea: **uncertainty compounds, it does not resolve.** A
% transcription that smooths over a doubtful word has destroyed the only thing
% it was made for — a reader must be able to tell, at every point, what was on
% the page and what was supplied. The apparatus macros below are therefore the
% critical apparatus, not decoration.
%
% The same commands are recognised by `scripts/render.mjs`, which produces the
% site's reading view, and by `scripts/tei.mjs`. A macro added here must be
% added there too, or rendering fails loudly — which is the wanted behaviour.
%
% Comments here are English, like the rest of the repository. The documents
% this preamble sets are French, and that is deliberate: see the skills.

\ProvidesPackage{germain}[2026/09/15 Transcriptions of mathematicians' manuscripts in Gallica]

\usepackage{amsmath,amssymb,amsthm}
\usepackage{mathrsfs}
% Kept from the parent preamble so the renderer's diagram support stays
% available; these manuscripts draw no commutative diagrams, and nothing here uses it.
\usepackage{tikz-cd}
\usepackage[english,french]{babel}
\usepackage{fontspec}

% --- Greek ------------------------------------------------------------------
% The text face stays fontspec's default, Latin Modern, which has no Greek:
% XeTeX drops every glyph it lacks with a line in the log and nothing on the
% page, and the Greek words the manuscripts quote vanished from the PDFs. So
% the Greek and Greek Extended blocks get a XeTeX character class of their own,
% and entering or leaving a run of it switches the family to CMU Serif — the
% same Computer Modern design, with polytonic Greek — and back. Only the
% family changes: a Greek word in \textit or \textbf keeps its shape and series.
% The transitions to and from every other class carry the tokens that class
% has towards an ordinary letter, so French spacing before « ; » or inside
% « » still holds after a Greek word. No grouping, so no brace or \endgroup
% can come out unbalanced; maths is left alone, since XeTeX inserts nothing
% there. Kept identical in `exercices.sty`.
\makeatletter
\newfontfamily\germainGreekfont{cmun}[
  Extension=.otf, NFSSFamily=germaingreek,
  UprightFont=*rm, ItalicFont=*ti, SlantedFont=*sl,
  BoldFont=*bx, BoldItalicFont=*bi, BoldSlantedFont=*bl]
\newXeTeXintercharclass\germain@greek
\def\germain@greekrange#1#2{%
  \@tempcnta=#1\relax
  \loop
    \XeTeXcharclass\@tempcnta=\germain@greek
    \ifnum\@tempcnta<#2\advance\@tempcnta\@ne\repeat}
\germain@greekrange{"0370}{"03FF}
\germain@greekrange{"1F00}{"1FFF}
\def\germain@greekprev{\rmdefault}
\def\germain@greekin{%
  \edef\germain@greekprev{\f@family}\fontfamily{germaingreek}\selectfont}
\def\germain@greekout{\fontfamily{\germain@greekprev}\selectfont}
\def\germain@greekjoin{%
  \XeTeXinterchartokenstate=\@ne
  \@tempcnta=\z@
  \loop
    \ifnum\@tempcnta=\germain@greek\else
      \XeTeXinterchartoks\germain@greek\@tempcnta=\expandafter{%
        \expandafter\germain@greekout\the\XeTeXinterchartoks\z@\@tempcnta}%
      \XeTeXinterchartoks\@tempcnta\germain@greek=\expandafter{%
        \the\XeTeXinterchartoks\@tempcnta\z@\germain@greekin}%
    \fi
    \ifnum\@tempcnta<4095\advance\@tempcnta\@ne\repeat}
% babel-french sets its own classes, and saves and restores the interchar
% state, each time a language is selected: join again after it does.
\addto\extrasfrench{\germain@greekjoin}
\addto\extrasenglish{\germain@greekjoin}
\AtBeginDocument{\germain@greekjoin}
\makeatother

\usepackage{geometry}
\geometry{a4paper,margin=25mm,marginparwidth=28mm}
\usepackage{marginnote}
\usepackage{xcolor}
\usepackage[normalem]{ulem}
\setlength{\emergencystretch}{3em}
\usepackage{hyperref}
\hypersetup{colorlinks=true,linkcolor=black,urlcolor=black,pdfborder={0 0 0}}

% --- Batch metadata ---------------------------------------------------------
% Taken as they stand by the rendering script, which shows them at the head of
% the reading view. They fix what the file claims to cover. The macro names are
% the parent project's, so that the scripts stay shared; \folder here means the
% volume's slug — `fr-9115` — and \pages counts Gallica views.

\newcommand{\folder}[1]{\gdef\grothFolder{#1}}
\newcommand{\batch}[1]{\gdef\grothBatch{#1}}
\newcommand{\pages}[2]{\gdef\grothFirst{#1}\gdef\grothLast{#2}}
\newcommand{\foldertitle}[1]{\gdef\grothFoldertitle{#1}}

% The holder's own shelfmark, printed as they write it: « Français 9115 ».
\newcommand{\shelfmark}[1]{\gdef\grothShelfmark{#1}}

% The dating, copied verbatim from the holder's catalogue — never inferred
% here. « XIXe siècle » is the BnF's; a pass that dates a leaf from its content
% says so in a \note{}, as its own inference.
\newcommand{\dating}[1]{\gdef\grothDating{#1}}

% The status notice, printed in the title block and repeated as a diagonal
% watermark on every page of the PDF. The manuscripts are in the public
% domain; what this line declares is not a right but a quality — an
% unverified first pass by a machine. The skills write the line; a file
% without it gets no watermark, so leaving it out is visible in the output.
\newcommand{\watermark}[1]{\gdef\grothWatermark{#1}}

\gdef\grothFolder{?}\gdef\grothBatch{}\gdef\grothFirst{?}\gdef\grothLast{?}%
\gdef\grothFoldertitle{}\gdef\grothShelfmark{}\gdef\grothDating{}\gdef\grothWatermark{}

% --- The résumé -------------------------------------------------------------
\newenvironment{resume}
  {\small\setlength{\parskip}{0.45em}}
  {\par\medskip\hrule\medskip}

% --- Keywords ---------------------------------------------------------------
% In English, closing the modernised reading's résumé: the modern vocabulary
% under which to search for what the leaves construct. The manifest extracts
% them to tag the volume — this line is the single source of the tags.
\newcommand{\keywords}[1]{\par\smallskip\noindent
  {\footnotesize\textsc{Keywords} --- \textit{#1}}\par}

% --- The critical apparatus -------------------------------------------------

% The Gallica view number, in the margin. This is the anchor that lets a
% disagreement over a formula be settled by looking at *one* image rather than
% a volume of 750. The site uses it to turn the facsimile as the transcription
% is scrolled. A view is one scanned image: usually one side of a leaf.
\newcommand{\page}[1]{%
  \marginnote{\footnotesize\color{gray!70}\textsf{v.\,#1}}%
  \label{page:#1}\ignorespaces}

% The BnF's pencilled foliation, where the leaf carries it and a pass can read
% it: « 198r ». This is what the literature cites — « ff. 198r–208v » — and
% the one line that turns a citation into a view. Recorded, never inferred:
% a folio a pass cannot read is not written.
\newcommand{\folio}[1]{%
  \marginnote{\footnotesize\color{gray!50}\textsf{f.\,#1}}\ignorespaces}

% The modernised reading's coarser counterpart to \page{}: which views a
% section is drawn from.
\newcommand{\pagerange}[2]{%
  \marginnote{\footnotesize\color{gray!70}\textsf{v.\,#1--#2}}%
  \ignorespaces}

% Illegible: never guessed.
\newcommand{\ill}{\textcolor{red!60!black}{[\,\ldots\,]}}

% A doubtful reading: the word is given, and flagged as uncertain.
\newcommand{\uncertain}[1]{\uline{#1}}

% An editorial addition — a missing letter, an implied word.
\newcommand{\add}[1]{\textcolor{blue!50!black}{[#1]}}

% Struck out by the writer. What was crossed out is part of the document.
\newcommand{\struck}[1]{\sout{#1}}

% The transcriber's note, on the state of the page or the mathematics.
\newcommand{\note}[1]{\par\smallskip{\small\color{gray!60!black}\textit{[#1]}}\par\smallskip}

% A marginal note of hers, distinct from the body of the text.
\newcommand{\marginal}[1]{\par\smallskip\noindent\hspace{1em}%
  {\small\color{gray!60!black}#1}\par\smallskip}

% --- Title ------------------------------------------------------------------
% The title block is French, like the documents themselves.

\AddToHook{shipout/background}{%
  \ifx\grothWatermark\empty\else
    \begin{tikzpicture}[remember picture, overlay]
      \node[rotate=52, scale=3.4, align=center, text=gray!18,
            font=\sffamily\bfseries]
        at (current page.center) {\grothWatermark};
    \end{tikzpicture}%
  \fi}

\AtBeginDocument{%
  \begin{center}
    {\large\bfseries Manuscrits de mathématiciens (Gallica) ---
      \ifx\grothShelfmark\empty\grothFolder\else\grothShelfmark\fi}\\[2pt]
    {\itshape\grothFoldertitle}\\[4pt]
    {\small vues \grothFirst--\grothLast
      \ifx\grothBatch\empty\else\ (lot \grothBatch)\fi}\\[2pt]
    \ifx\grothDating\empty\else
      {\small\color{gray!60!black}Datation du catalogue : \grothDating}\\[2pt]
    \fi
    \ifx\grothWatermark\empty\else
      {\small\bfseries\color{red!45!gray}\grothWatermark}\\[2pt]
    \fi
    {\footnotesize\color{gray!60!black}Transcription établie d'après les images IIIF de Gallica.
     Source gallica.bnf.fr / Bibliothèque nationale de France.\\
     Les lectures incertaines sont soulignées, les illisibles notées [\,\ldots\,],
     les ajouts entre crochets ; \textsf{v.} numérote les vues de Gallica,
     \textsf{f.} la foliotation au crayon de la BnF.}
  \end{center}
  \vspace{1em}}

\endinput


\folder{fr-9115}
\shelfmark{Français 9115}
\batch{5}
\pages{81}{100}
\dating{1801-1900}
\watermark{Lecture automatique\\non vérifiée}
\foldertitle{Papiers de Sophie GERMAIN. — Recueil de dissertations et problèmes mathématiques et physiques.}

\begin{document}

\note{Numérotation des feuillets. La série à l'encre se poursuit sur les
rectos de la mise au net — vues 81, 83, 85, 87, 89 → 40, 41, 42, 43, 44 — avec
à côté, sur les trois premières, la pagination de l'auteur, impaire et biffée,
dont 39 seul a pu être lu (vue 83) ; la vue 87 ne porte qu'un « 43 », non
biffé. La pagination paire est en haut à gauche de chaque verso, non biffée :
38, 40, 42, 44 aux vues 82, 84, 86, 88. Les feuillets de travail portent la
même série et rien d'autre : 45 (vue 91), 46 (vue 93), 47 (vue 95, avec la
lettre « E »), 48 (vue 97), 49 (vue 99, au recto blanc du feuillet écrit à la
vue 100). Tous ces nombres ont été lus sur la vue ; aucun n'est déduit. Ils
sont à l'encre : aucune foliotation au crayon de la Bibliothèque n'a été lue
sur ces vues, et aucun « f. » n'est donc porté ici. Les vues 90, 92, 94, 96,
98 et 99 sont blanches et ne sont pas transcrites.}

\page{81}
distances moyennes et la loi de la répartition de la courbure autour
d'un point choisi sur la surface.

Il faut d'abord se rappeler qu'$AC$ est le développement des bases des
deux \emph{surfaces des distances} ; que ces bases sont contenues dans
le plan tangent ; et que, par conséquent $AC$ représente la position
de ce plan.

Cela posé, on verra aisément que, pour toutes les figures de la
surface comprises entre la forme cylindrique et celle de la surface
dont les rayons de courbures principales, et, par conséquent, aussi
les \emph{lignes de distances} principales sont égaux et de signes
contraires, les angles en $q$ et $q'$ ; $h$ et $h'$ ; demeurent les
mêmes (fig. 2, 3 et 4).

La valeur de ces angles dépend, en effet, uniquement et, de
l'intervalle $pq$, invariable, puisque, par construction, il est, dans
tous les cas possibles, la quatrième partie de $AC$, et de la somme
des deux lignes $pq$, $gh$ ; qui, ainsi qu'on vient de le dire ; est
également invariable.

Lorsque les deux rayons de courbures principales sont dirigés en sens
opposés, la ligne brisée $Bqh$ \&c, et dès lors aussi la directrice de
la \emph{surface des distances} dont cette ligne est le développement,
est donc indépendante du rapport entre les deux rayons de courbures ;
et, par conséquent, elle est, dans tous les cas, celle qui convient à
la surface cylindrique fig. 2, limite commune entre ce genre de
surfaces et celles dont les rayons de courbures sont dirigés du même
côté. Les différences entre ces diverses surfaces consistent seulement
en ce que des parties, plus ou moins grandes de la directrice se
trouvent placées au dessus du plan tangent.

On se représentera les conséquences de la transformation qui feroit
passer la surface, de la forme cylindrique à celle où les rayons

\page{82}
de principales courbures sont égaux mais de signes contraires ; si on
conçoit que la ligne $AC$ fig. 2, se meuve parallèlement à elle-même
en se rapprochant graduellement de $BD$ fig. 3, jusqu'à ce que les deux
lignes se confondent fig. 4. En effet, durant ce mouvement, la
\emph{ligne $pq$ de principale distance} située au dessus du plan
tangent, ligne qui étoit nulle à l'égard de la surface cylindrique,
augmentera par degrés, en même tems que $gh$ diminuera de quantités
égales, jusqu'à ce qu'enfin ces deux lignes, terminées aux côtés
opposés de $AC$, deviennent égales.

On a vu dans les N°s précédens qu'en passant de la
surface sphérique à la surface cylindrique, par le nombre infini des
configurations qui les séparent, la ligne brisée $Bqh$ \&c, et, par
conséquent, aussi la directrice de la \emph{surface des distances},
dont cette ligne est le développement, change de formes autant de
fois. On a vu en même tems qu'au milieu de tous ces changemens, la
\emph{surface des moyennes distances} n'en subit aucun.

Il résulte de ce qu'on vient de dire que le contraire arrive quand on
passe de la surface cylindrique à celle dont les rayons de courbures
principales sont égaux, mais de signes contraires. La directrice de la
\emph{surface des distances} conserve la même forme ; mais une partie
de cette courbe passe de l'autre côté du plan tangent, et, par
conséquent, la \emph{surface des distances moyennes} diminue
graduellement, jusqu'à ce qu'elle s'anéantisse, lorsque les
\emph{lignes de distances} terminées à une des faces du plan tangent
sont égales aux lignes de distances terminées à l'autre face du même
plan.

Il nous reste à montrer que, malgré les différences

\page{83}
qu'entraîne le changement de direction d'un des rayons de courbures
principales, la loi de la distribution de la courbure autour d'un des
points de la surface ne souffre aucune altération ; et ainsi
\uncertain{il en est} encore vrai, à l'égard des surfaces dont les
rayons de principales courbures sont dirigés en sens contraires :

que la somme des \emph{lignes de distances} contenues dans deux plans
normaux perpendiculaires entr'eux est indépendante du choix de ces
plans :

que la \emph{ligne de distance} contenue dans le plan normal qui fait
l'angle de $45^\circ$ avec ceux des courbures principales est moyenne
entre les \emph{lignes de distances principales} :

Enfin qu'en embrassant le contour entier du point donné, il y a, en
général, quatre directions du plan normal pour lesquelles les
\emph{lignes de distances} sont égales, et que le nombre de ces
directions se réduit à deux à l'égard des \emph{lignes de distances
principales}.

Ces diverses propositions ayant été établies dans le cas où les deux
rayons de principales courbures sont dirigés dans le même sens, il
suffit, pour les étendre au cas présent, d'observer que,
l'augmentation des \emph{lignes de distances} situées au dessus de
$AC$ étant considérée, à raison du signe, comme une diminution, les
diverses \emph{lignes des distances} dont la comparaison a servi à
démontrer, dans le cas de la surface cylindrique, les propositions
dont il s'agit, éprouvent à la fois des diminutions proportionnelles ;
en sorte que, leur grandeur relative ne changeant pas, les mêmes
propositions appuyées sur des raisonnemens semblables demeurent dans
leur entier.

\page{84}
En parlant de la surface cylindrique, on a déjà remarqué que, deux
points, seulement, du développement de la directrice des surfaces
\emph{des distances} savoir $q$ et $q'$, tombant sur $AC$, une des
courbures principales est nulle. Quand les deux rayons de courbures
principales sont de signes opposés, quatre points du développement de
la directrice, savoir $v$ et $t$ ; $v'$ et $t'$ tombant sur $AC$, ni
l'une ni l'autre des courbures principales ne sont nulles ; mais,
comme, à l'égard des surfaces cylindriques, la sphère dont le rayon
est infini est au nombre de celles qui ont leur centre sur la normale
au point donné

Dans le cas où les deux rayons de courbures principales sont égaux et
de signes contraires, les points en $v$ et $u$ ; $v'$ et $u'$ se
confondent avec ceux en $n$, $n'$ ; et $n''$, $n'''$ lesquels
partagent $BD$ en quatre parties égales. La sphère dont le rayon est
infini est donc alors celle de moyenne courbure, et cette sphère se
confond avec le plan tangent.
\note{les quatre points sont nommés $v$, $t$, $v'$, $t'$ au premier
alinéa et $v$, $u$, $v'$, $u'$ au second : les deux lectures ont été
faites sur la vue et ne sont pas harmonisées ici. Le premier $n$ est
écrit sur une lettre reprise à la plume.}

N° 9 Si la considération des courbures moyennes mérite quelque
attention c'est surtout à cause du jour qu'elle doit répandre sur les
questions où la courbure des surface joue le rôle d'une quantité
dynamique. Peut-être même seroit-il impossible de se rendre un compte
satisfaisant d'une question de ce genre ; embrassée dans toute sa
généralité, si on ne connoissoit ni la loi de la répartition de la
courbure autour d'un point donné, ni l'existence des courbures
moyennes qui se lie à celle de cette loi.

Supposons que les forces qui agissent sur un des points de la surface
sont proportionnelles à sa courbure. On sait que, quelles que soient
le nombre et la direction des forces qui agissant sur un point, elles
peuvent être séparées en
\note{« qui agissant sur un point » est bien ce que porte le feuillet.}

\page{85}
deux parties qui agiront suivant deux plans perpendiculaires entr'eux
menés par ce point : on sait encore que la position des plans choisis
est arbitraire.

Tant qu'on ne distingue pas l'idée de la courbure de celle de la
figure des surfaces, une telle séparation et l'équivalence d'un nombre
infini de positions différentes des plans choisis doivent paroître
inconcevables. Comment imaginer, en effet, que la figure d'une
surface, séparée en deux parties, se réduise à deux courbes linéaires,
et que ces courbes puissent être remplacées par une infinité d'autres
courbes de figures entièrement différentes.

La notion des courbures moyennes et la loi de la distribution de la
courbure, telle qu'elle a été expliquée ; satisfont, au contraire à
toutes les conditions qu'exigent les principes de la mécanique.

Si les forces qui agissent sur un des points de la surface sont
proportionnelles à la courbure, elles sont réparties dans tout le
contour de ce point ; car la courbure elle même est distribuée dans le
contour entier. Le nombre de ces forces est infini ; d'elles
\struck{\ill{}} chacune infiniment petite\struck{s}, \struck{et} elle
est proportionnelle à la ligne de distance contenue dans le plan
\struck{dela} sa direction.
\note{« chacune » est écrit au-dessus de la ligne, sur un mot biffé
d'un trait épais et illisible ; le « s » de « petites », le mot « et »
et le mot « dela » sont biffés de la même encre. La phrase se lit donc
sur le feuillet : « Le nombre de ces forces est infini ; d'elles
chacune infiniment petite, elle est proportionnelle à la ligne de
distance contenue dans le plan sa direction ».}

L'action exercée doit être indépendante de la manière dont les forces
qui l'exercent sont réparties autour du point donné ; quelles que
soient les différences, cette action sera donc toujours égale à celle
qui seroit due à une répartition uniforme des mêmes forces.

Nous avons vû aussi que la courbure moyenne de la surface est
indépendante de la forme de la \emph{surface des distances},

\page{86}
surface qui représente, dans le cas présent, la distribution des
forces dues à la courbure : l'action exercée est proportionnelle à la
\emph{surface des distances moyennes} ; cette action seroit donc
encore la même si le point donné appartenoit, non plus à la surface
que l'on considère, mais à la sphère de moyenne courbure de cette
surface ; c'est à dire, si la distribution des forces autour du point
donné étoit uniforme.

Les forces dues à la courbure, aussi bien que tout autre genre de
forces qui agiroient sur un point donné, peuvent être décomposées
suivant deux plans perpendiculaires entr'eux dont la position,
d'ailleurs arbitraire, est assujettie à passer par le point donné.
Dans le cas présent, les résultantes des forces qui agissent sur le
point donné de la surface, seront proportionnelles aux lignes de
distances contenues dans les plans choisis ; et nous avons démontré
que la somme de ces deux lignes ne dépend pas non plus de la position
du système des plans qui les contiennent.

Il est bon de faire observer que les \emph{lignes de distances} sont
ici, proportionnelles à des quantités finies ; tandis que, dans ce qui
précède, elles sont regardées comme proportionnelles à des quantités
infiniment petites, savoir à chacune des forces qui agissent sur le
point donné. Ces suppositions ne sont pas contradictoires ; mais elles
résultent de deux points de vues différens d'un même objet. L'un est
conforme au principe général de la composition des forces : l'autre,
tout spécial, s'attache à chacune des forces qui agissent autour du
point donné de la surface ; et, parce que le nombre de ces forces est
infini, leur ensemble se trouve représenté par la \emph{surface des
distances moyennes}.
\note{le bord droit du feuillet passe sous le montage : les dernières
lettres de plusieurs lignes sont cachées, « agissent » et
« contradictoires » entre autres, dont le début seul se lit sur la
vue.}

\page{87}
L'emploi de la composition des forces mène à des formules commodes ;
mais, quand on cherche à se représenter l'action des forces dues à la
courbure, la considération des \emph{surfaces des distances} convient
mieux.

Et, en effet, les \emph{surfaces des distances} sont terminées, d'une
part au plan tangent et de l'autre à la surface donnée, sur laquelle
elles s'appuyent dans tous les points de leurs directrices. Si on
conçoit que ces surfaces, devenues solides, continuent d'être
interposées entre la surface donnée et son plan tangent, elles
pourront être regardées comme l'obstacle qui empêche les points des
directrices, c'est-à-dire, les points qui environnent celui de
tangence, de s'approcher de ce plan. Les différens points dont se
composent les \emph{surfaces des distances}, ainsi interposées, étant
doués de forces égales, l'étendue de ces surfaces représentera
fidèlement les forces dues à la distance entre les directrices et le
plan tangent, c'est à dire, qu'elle représentera l'action des forces
nées de la courbure de la surface donnée.

Il existe, comme on l'a vu, les plus grands rapports entre le principe
de la composition des forces et la question des courbures moyennes ;
et il suffit même d'avoir à considérer des forces dues à la courbure
d'une surface pour être conduit à \struck{distinguer} ne pas confondre
la figure, avec la courbure envisagée comme donnant naissance à une
quantité dynamique.
\note{« ne pas confondre » est écrit au-dessus de « distinguer »
biffé ; « avec » est écrit par-dessus un mot que la surcharge rend
illisible.}

La considération des courbures moyennes est nécessairement applicable
à toutes les questions où la courbure des surfaces joue le rôle d'une
puissance ; mais, le but que nous nous sommes proposé n'étant pas de
traiter ces questions ; nous nous bornerons donc, à l'égard des
applications, à un petit nombre de remarques.
\note{« n'étant » est écrit sur un mot repris à la plume, illisible
sous la correction.}

L'équation des surfaces élastiques ayant été l'occasion des recherches
qu'on vient de lire, il est à propos de faire observer d'abord, que,
pour justifier l'hypothèse qui \uncertain{a} servi à trouver, la
première fois,

\page{88}
l'équation de ces surfaces, il suffiroit de répéter les diverses
propositions qui ont été démontrées, en y joignant la désignation de ce
genre de surfaces.

La simplicité de description et le petit nombre des figures
différentes de la surface qui peuvent convenir à la construction des
voûtes ont sans doute masqué aux yeux des géomètres qui se sont
occupés de cette matière, la nécessité de savoir exprimer, d'une
manière générale, les résistances nées de la courbure des surfaces.

La notion des courbures moyennes ne sauroit cependant demeurer
étrangère à de telles recherches. Supposons, par exemple, qu'on
veuille comparer les voûtes en dômes aux voûtes en berceaux ; ne
sera-t-il pas de la plus grande utilité de savoir aussi comparer les
deux genres de courbures \uncertain{employés}

Nous terminerons ces réflexions en présentant l'exemple de deux
surfaces de figures différentes qui possèdent, en même tems qu'une
étendue égale, la même courbure moyenne dans tous leurs points.

Les deux surfaces dont nous voulons parler sont celles du cylindre
circonscrit et la quatrième partie de la surface sphérique décrite du
rayon double de celui de la sphère inscrite. Ces deux surfaces sont
d'une étendue égale ; puisque le cylindre et la sphère inscrite sont
égaux en surface, et que celle des sphères est en raison du quarré de
leurs rayons. De plus, ces mêmes surfaces ont dans tous leurs points
la même courbure moyenne ; car, d'après ce qui précède, la sphère dont
le rayon est double de celui de la base de la surface cylindrique,
n'est autre chose que la sphère \emph{de moyenne courbure} de cette
surface
\note{la fin de cette ligne passe sous le montage, comme à la vue 86 :
de « employés » les dernières lettres sont cachées, et la ponctuation
qui fermoit l'interrogation n'est pas lue. « de moyenne courbure »,
à la dernière ligne du feuillet, est souligné deux fois.}

\page{89}
\note{planche des figures du mémoire, à l'encre, sans une ligne de
texte ; le nombre 44 est en haut à droite. Quatre figures y sont
tracées l'une sous l'autre et numérotées « fig. 1 » à « fig. 4 » dans
la marge de droite. Chacune est un parallélogramme $ABDC$ — $A$ et $C$
aux extrémités de la ligne supérieure, $B$ et $D$ de l'inférieure —
que traverse une ligne brisée en dents, avec les perpendiculaires
ponctuées ou pleines qui joignent les deux lignes. La ligne
supérieure porte, de $A$ vers $C$ : $p$, $m$, $q$, $m'$, $p'$, $m''$,
$q'$ ; l'inférieure, de $B$ vers $D$ : $h$, $n$, $h'$, $n''$. Les
sommets de la brisée sont nommés $q$, $q'$ en haut et $h$, $h'$ en
bas, et la figure 1 porte en outre $q$ et $t$ joints par une ligne
ponctuée. La figure 3 ajoute sur la ligne supérieure les points $v$,
$u$ et $v'$, $u'$ de part et d'autre de $p$ et de $p'$, et la figure 4
ramène $A$ sur $B$ et $C$ sur $D$, les deux lignes du parallélogramme
étant alors confondues. Ce sont les figures que le mémoire appelle
depuis la vue 73 ; aucune n'est légendée autrement que par son
numéro.}

\page{91}
Moyennes courbures
\note{le feuillet ne porte que ce titre, tracé d'une grande écriture au
milieu de la page et souligné ; le nombre 45 est à l'encre en haut à
droite. Le verso, vue 92, est blanc.}

\page{93}
Prenons la voûte en dôme formée par la révolution d'une demie courbe
quelconque. Les élémens compris entre deux plans parallèles aux pieds
de la voûte formeront une zone où la courbure sera partout semblable.
Les rayons de moyenne courbure de tous ces élémens seront en même tems
ceux des élémens d'une surface conique circulaire dont le rayon de la
base sera égal à la moitié du rayon de moyenne courbure.
\uncertain{Aussi on formera} ainsi diverses portions de surfaces
cylindriques dont la courbure sera la même que celle des diverses
zones de la surface de révolution.

Si on \uncertain{prenoit} de $A$ à $B$ \struck{de $A$ à $B$} tous les
rayons d'une \struck{1\textsuperscript{re}} zone et de $A'$ à $B'$
ceux d'une autre zone et ainsi de suite en avançant vers $C$, ceux des
diverses zones se succéderoient de $A$ \uncertain{à} $C$, si par les
extrémités de ces rayons pris comme ordonnées on feroit passer une
ligne droite cheminant de $A$ en $C$ elle s'appuyeroit sur toutes les
valeurs égales de $A$ en $B$ de $A'$ en $B'$ de \&c
\struck{\ill{}} engendreroit une surface cylindrique dont les élémens
auroient aussi la même courbure moyenne que ceux de la surface
\note{feuillet de travail, d'une écriture rapide et fort corrigée, et
déchiré : tout le bas du feuillet manque, la déchirure emportant la
suite du texte, dont les dernières lignes se serrent dans le coin qui
subsiste. Les couples de lettres sont écrits sur deux segments tracés
à la plume, $A$——$B$ au-dessus et $A'$——$B'$ au-dessous ; une grande
accolade verticale, tracée d'un seul trait, embrasse « ainsi de suite
en avançant » et la ligne suivante. Un « \&c » et le mot qui le suit
sont biffés d'une croix.}

\page{95}
N° \struck{\ill{}}

L'équation des surfaces étant définie il \struck{nous est à chercher}
\ill{} supposerons qu'une cause extérieure ait changé la forme du
suivant ce qui a été établi nous dirons que la \struck{surface}
question des surface consiste à exprimer quels sont les espaces
parcourus par les points d'un solide \struck{qui dans le changement de
forme que le solide} qui, dans le changement de forme qu'il a subi ce
solide a subi, \struck{ont été mus dans des directions perpendiculaires
aux plans}

Commencement de la p. A

\struck{Reprenons la} \struck{afin de mieux apprécier les parties
\uncertain{simplificatrices} qui caractérise le cas des surfaces nous
allons rappeler les principes généraux qui conviennent à ce genre de
questions}

Revenons maintenant présentement à la question générale

À $B$. Par rapport au point de tangence l'espace parcouru est
immédiatement mesuré par le changement de position du plan tangent
qui, durant le mouvement a été transporté parallèlement à lui-même
sans cesser de toucher la surface. On vient de dire que les points
situés dans la profondeur du solide subissent des déplacemens égaux à
ceux des points extérieurs cette circonstance permet de
\uncertain{substituer} la considération de l'élément de la surface à
celui du solide et il ne nous reste donc plus à déterminer que les
espaces parcourus par les points
\note{feuillet de travail, écrit d'une plume rapide et repris partout.
En haut à droite la lettre « E » suivie d'un signe empâté, et le
nombre 47. Une grande croix couvre les six premières lignes, qui sont
en outre biffées ligne à ligne ; les trois mots qui suivent « N° » ne
se laissent pas lire sous la rature. Les mots « à celui du solide »,
« ne », « donc plus » et « que » sont ajoutés entre les lignes de la
dernière phrase, et sont ici mis à leur place ; « généraux » est de
même ajouté au-dessus de la ligne dans le membre biffé. « Commencement
de la p. A » est une indication de l'auteur à elle-même, portée seule
sur sa ligne.}

\page{97}
Il est peut être inutile de remarquer que la figure nodale qui
manifeste la suite des points qui demeurent en repos pendant le
mouvement \ill{} un symptôme irrécusable de la figure complette
\ill{} l'invariabilité de l'épaisseur \ill{} substituer

la considération de l'élément d'un solide celle de l'élément de la
\struck{surface} superficie qui le recouvre et il suffit de
déterminer le déplacement des points superficiels pour connoître celui
des points situés dans la profondeur du solide \struck{cette} ces
conclusion résulte évidemment de l'uniformité en même tems des espaces
parcourus, par les points qui composent les couches intérieures du
solide, et
\marginal{cette uniformité elle même résulte de l'invariabilité de
l'épaisseur}
dans le cas des surfaces donc l'épaisseur \struck{qui ne peut influer
\ill{} sur} la figure qu'affectent les faces du solide \struck{elle}
influe uniquement sur la vitesse du mouvement, \struck{les
différences} l'élasticité naturelle influe \struck{aussi} uniquement
sur cette \struck{même} vitesse \struck{et par conséquent} il suffit
alors de multiplier l'équation trouvée pour la surface par le
coefficient constant \struck{qui} composé des deux facteurs qui
expriment l'influence de l'épaisseur et celle de l'élasticité
naturelle pour étendre l'équation qui a été formée sans qu'on ait rien
statué sur \add{à} l'égard de ces deux qualités du solide à ce même
solide considéré dans la réalité de son existence physique.
\struck{\ill{}}

Si cette manière d'envisager la question pouvoit encore laisser
quelque doute dans l'esprit du lecteur j'ajouterois \struck{\ill{}}
qu'il seroit \struck{et} tout à fait impossible de concevoir de se
rendre compte de la manière dont l'équation des plaques élastiques a
été trouvée
\note{feuillet de travail, numéroté 48 dans la marge de droite, repris
de bout en bout : mots biffés, remplacemens entre les lignes, une
addition dans la marge de gauche. Les quatre lignes du haut sont
écrites après coup, d'une plume très fine et très serrée, dans
l'espace laissé libre au-dessus du premier alinéa ; elles ne se
laissent lire qu'en partie, et ce qui manque y est marqué comme partout
ailleurs. Un
grand crochet embrasse, à gauche, les six lignes qui vont de « il
suffit alors de multiplier » à « l'élasticité naturelle » ; les mots
« alors », « superficie », « surface », « en même tems »,
« uniquement » et « de concevoir » sont ajoutés au-dessus de la
ligne, et sont ici mis à leur place. Le chiffre « 1 » est porté dans
la marge de gauche en regard de « j'ajouterois ».}

\page{100}
\marginal{sur la surface}
qui environnent le point de tangence \struck{et qui} par rapport à
chacun \struck{Il est évident que les espaces parcourus seront connus
si on sait à la fois qu'il est, dans l'état de repos, la distance
d'un point voisin au plan tangent, quel déplacement le plan tangent a
subi pendant le mouvement et enfin durant le mouvement et}
\marginal{par rapport à chacun de ces points}

On sait déjà quel a été le changement de position du plan tangent il
est donc évident que les espaces parcourus par les points qui
environnent le point de tangence seront connus si on sait à la fois
par rapport à chacun de ces points quel est sa distance au plan
tangent, dans l'état naturel de la surface et \struck{quelle est elle}
et après qu'une cause extérieure l'a forcée de changer de forme.

Les distances dont nous parlons sont liées entr'elles par une loi
facile à reconnaître cette loi constitue la notion des courbures
moyennes nous allons l'expliquer

\ill{}
\note{feuillet de travail. Le premier alinéa est barré d'un grand
trait diagonal, outre les ratures ligne à ligne ; « parcourus » et
« par rapport à chacun de ces points » sont ajoutés au-dessus des
lignes. La dernière ligne, sous « nous allons l'expliquer », est
réduite à un mot et à des taches d'encre, et n'est pas lue. Le recto
de ce feuillet, vue 99, ne porte que le nombre 49.}

\end{document}
